\subsection{A note on the predictions of models with modular flavor symmetries --- arXiv:1909.06910}

\textbf{Authors:} Mu-Chun Chen, Sa{\'u}l Ramos-S{\'a}nchez, Michael Ratz

\paragraph{主な主張}
モジュラーフレーバー対称性だけではK\"ahlerポテンシャルは固定されず、許される非正則項に含まれる追加パラメータが質量・混合の予言力を低下させると指摘した\cite{Chen:2019ewa}。

\paragraph{新規性}
正則な超ポテンシャルだけを用いた予言に対し、運動項の正準規格化による補正を明示し、予言力維持に必要な追加構造を論じた。

\paragraph{位置付け}
モジュラーフレーバー模型の予言がどの仮定に依存するかを検討する際の重要な注意点を与える。

\paragraph{コメント（読書メモ）}
読書メモ：このK\"ahlerポテンシャル由来の予言力の問題は「無視され続けている」と感じた（2024年10月の読書時点での私見であり、分野全体について確認した評価ではない）。
